\section{Resultados da recuperação textual}

O arquivo docling-langchain-bge-m3-retrieval-pages apresenta o resultado da execução bem sucedida da pipeline, essencial para inspeção manual dos resultados e construção de visualizações, onde é possível encontrar os chunks ranqueados e as métricas calculadas por pergunta.

A saída legível contém, para cada pergunta, informações sobre o documento avaliado, a pergunta utilizada como consulta, as páginas esperadas como evidência e os resultados recuperados pelo pipeline. No exemplo analisado, referente ao documento PH\_2016.06.08\_Economy-Final.pdf, a página de evidência esperada era a página 5. O pipeline recuperou chunks textuais associados às páginas 5 e 6, sendo que o primeiro resultado ranqueado já continha a página correta. Dessa forma, a pergunta obteve first\_hit\_rank igual a 1 e MRR igual a 1,0. O "score" que é visto no ranking é referente à nota de similaridade com a consulta, logo, quanto maior o número, maior a coesão daquele chunk ou página com o contexto.


\begin{figure}[h!]
    \centering
    \includegraphics[width=1.0\textwidth]{02_Textuais/desenvolvimento/metodologia/imagens/JSON 1.png}
    \caption{Arquivo resultado do pipeline Docling}
    \label{fig:minha_imagem}
\end{figure}

\begin{figure}[h!]
    \centering
    \begin{minipage}{0.48\textwidth}
        \centering
        \includegraphics[width=\textwidth]{02_Textuais/desenvolvimento/metodologia/imagens/Captura de Tela 2026-06-10 às 16.48.26.png}
        \caption*{(a) Exemplo do ranking de chunks}
    \end{minipage}
    \hfill
    \begin{minipage}{0.48\textwidth}
        \centering
        \includegraphics[width=\textwidth]{02_Textuais/desenvolvimento/metodologia/imagens/Captura de Tela 2026-06-10 às 16.49.06.png}
        \caption*{(b) Exemplo do ranking de páginas}
    \end{minipage}
    \label{fig:mmlongbench-doc-visao-geral}
\end{figure}


Os resultados recuperados aparecem principalmente em duas estruturas: ranked\_chunks e ranked\_pages. A primeira lista os chunks textuais retornados, ordenados pela pontuação de similaridade com a pergunta. Cada chunk possui uma posição no ranking, um identificador, uma pontuação, uma ou mais páginas associadas e a indicação se ele coincide com alguma página de evidência. A segunda estrutura, ranked\_pages, consolida esses resultados no nível de página, métrica que é o padrão para o estudo, removendo repetições e indicando em que posição do ranking a página apareceu pela primeira vez.

A~\ref{tab:saida-textual-readable} resume os principais campos presentes na saída legível do pipeline textual.

\begin{table}[htbp]
\centering
\caption{Principais campos da saída legível do pipeline textual.}
\label{tab:saida-textual-readable}
\small
\renewcommand{\arraystretch}{1.2}
\begin{tabular}{|p{4.2cm}|p{9.0cm}|}
\hline
\textbf{Campo} & \textbf{Descrição} \\
\hline
arch & Identifica a arquitetura utilizada no experimento, neste caso Docling, LangChain, BGE-M3 e Milvus. \\
\hline
model\_name & Indica o modelo de embeddings utilizado para representar as perguntas e os chunks textuais. \\
\hline
doc\_id & Identifica o PDF associado à pergunta avaliada. \\
\hline
question & Pergunta utilizada como consulta na etapa de recuperação. \\
\hline
evidence\_pages & Lista de páginas esperadas como evidência, de acordo com o MMLongBench-Doc. \\
\hline
parser & Indica o modo de processamento documental utilizado pelo Docling. \\
\hline
chunk\_token\_size & Tamanho máximo adotado para os chunks textuais. \\
\hline
page\_provenance & Indica a origem da informação de página usada para mapear chunks recuperados para páginas do PDF. \\
\hline
\end{tabular}
\end{table}


Por fim, o campo metrics apresenta as métricas calculadas para a pergunta. No exemplo apresentado, a primeira página relevante aparece na primeira posição do ranking, resultando em first\_hit\_rank igual a 1 e MRR igual a 1,0. Como a única página de evidência esperada foi recuperada já no primeiro resultado, os valores de Hit@1, Hit@3, Hit@5, Hit@10 e Hit@20 são verdadeiros, enquanto Evidence Page Recall@1, Recall@3, Recall@5, Recall@10 e Recall@20 assumem valor 1,0. Esse exemplo representa um caso em que o pipeline textual conseguiu recuperar corretamente a evidência esperada já no topo do ranking.


\begin{table}[htbp]
\centering
\caption{Interpretação dos resultados recuperados pelo pipeline textual.}
\label{tab:rankings-pipeline-textual}
\small
\begin{tabular}{p{4.2cm}p{9.0cm}}
\toprule
\textbf{Campo} & \textbf{Descrição} \\
\midrule
\texttt{ranked\_chunks} & Lista de chunks textuais recuperados, ordenados por similaridade com a pergunta. \\
\addlinespace
\texttt{rank} & Posição do chunk no ranking de recuperação. \\
\addlinespace
\texttt{score} & Pontuação de similaridade atribuída ao chunk recuperado. \\
\addlinespace
\texttt{page\_numbers} & Páginas de origem associadas ao chunk recuperado. \\
\addlinespace
\texttt{matches\_evidence} & Indica se o chunk recuperado está associado a alguma página esperada como evidência. \\
\addlinespace
\texttt{ranked\_pages} & Lista consolidada de páginas recuperadas a partir dos chunks. \\
\addlinespace
\texttt{first\_chunk\_rank} & Primeira posição em que uma página apareceu por meio de algum chunk recuperado. \\
\bottomrule
\end{tabular}
\end{table}


\subsection{Tempo de processamento da recuperação textual}

Previamente à análise de métricas de desempenho, é interessante analisar o tempo de processamento do pipeline, visto que o custo computacional representa uma métrica relevante ao trabalho quando comparada com as futuras arquiteruras que apresentam maior complexidade operacional, em especial a híbrida.

A~\ref{tab:tempo-processamento-textual} apresenta os principais tempos registrados no processamento dos documentos. Ao todo, foram processados 3.326 chunks textuais, onde não houveram falhas de execução.

\begin{table}[htbp]
\centering
\caption{Tempo de processamento do pipeline textual.}
\label{tab:tempo-processamento-textual}
\small
\begin{tabular}{p{7.8cm}p{4.8cm}}
\toprule
\textbf{Etapa ou métrica} & \textbf{Valor} \\
\midrule
Total de PDFs processados & 30 \\
Total de perguntas avaliadas & 172 \\
Total de chunks textuais gerados & 3.326 \\
Falhas de execução & 0 \\
\addlinespace
Parsing com Docling & 58min 25s \\
Indexação no Milvus & 2min 03s \\
Processamento textual total & 1h 00min 28s \\
Retrieval das perguntas & 11,88 s \\
Tempo total incluindo retrieval & 1h 00min 40s \\
\addlinespace
Média de processamento por PDF & 2min 01s \\
Mediana de processamento por PDF & 1min 10s \\
Média de retrieval por pergunta & 0,07 s \\
Mediana de retrieval por pergunta & 0,06 s \\
\bottomrule
\end{tabular}
\end{table}

Os resultados indicam que a etapa dominante do pipeline textual foi o parsing com Docling, responsável pela maior parte do tempo total de execução. Em comparação, a indexação vetorial com BGE-M3 e Milvus Lite apresentou custo relativamente baixo, consumindo pouco mais de dois minutos no conjunto completo. Após a construção dos índices, a etapa de retrieval foi rápida, totalizando 11,88 segundos para as 172 perguntas avaliadas.

Também foi observada variação significativa entre os documentos. O PDF mais rápido levou 11,98 segundos para ser processado, enquanto o mais lento levou 951,21 segundos, aproximadamente 15min 51s. Essa diferença indica que o custo do pipeline textual depende fortemente das características internas de cada PDF, como tamanho, estrutura, complexidade de layout e dificuldade de extração pelo parser, detalhe que então visa ser explorado na seção seguinte.

\subsection{Resultados detalhados por PDF}

Precedentemente à análise global, é válido ressaltar as métricas examinadas por documento. Essa análise permite observar variações de desempenho entre os PDFs avaliados, considerando que os documentos possuem tamanhos, estruturas e tipos de evidência diferentes. Nas tabelas a seguir, Qs indica a quantidade de perguntas associadas ao PDF, enquanto Falhas representa perguntas que não puderam ser avaliadas. A coluna FHM corresponde ao FirstHit médio, isto é, à posição média da primeira evidência correta entre as perguntas em que houve ao menos um acerto.


\begin{table}[htbp]
    \centering
    \caption{Tempo de processamento por PDF do pipeline textual.}
    \label{tab:tempo-processamento-textual-pdf}
    \scriptsize
    \resizebox{\textwidth}{!}{
    \begin{tabular}{r r r r r r r}
    \toprule
    \textbf{PDF} & \textbf{Qs} & \textbf{Chunks} & \textbf{Parsing (s)} & \textbf{Index. (s)} & \textbf{Proc. total (s)} & \textbf{Retrieval (s)} \\
    \midrule
    01 & 7 & 54 & 20.20 & 2.49 & 22.69 & 0.44 \\
    02 & 4 & 34 & 17.33 & 2.59 & 19.92 & 0.16 \\
    03 & 7 & 5 & 79.89 & 0.81 & 80.70 & 0.54 \\
    04 & 7 & 61 & 76.70 & 2.76 & 79.46 & 0.82 \\
    05 & 6 & 94 & 947.43 & 3.78 & 951.21 & 0.61 \\
    06 & 7 & 24 & 200.53 & 1.38 & 201.91 & 0.42 \\
    07 & 7 & 28 & 164.97 & 1.30 & 166.27 & 0.18 \\
    08 & 8 & 38 & 127.57 & 1.20 & 128.77 & 0.43 \\
    09 & 6 & 130 & 13.39 & 6.26 & 19.65 & 0.34 \\
    10 & 6 & 133 & 36.19 & 5.57 & 41.76 & 0.39 \\
    11 & 6 & 196 & 51.09 & 8.58 & 59.67 & 0.43 \\
    12 & 5 & 152 & 30.51 & 3.95 & 34.46 & 0.32 \\
    13 & 4 & 170 & 26.50 & 4.25 & 30.75 & 0.26 \\
    14 & 4 & 91 & 20.94 & 2.89 & 23.83 & 0.33 \\
    15 & 5 & 192 & 76.55 & 4.88 & 81.43 & 0.34 \\
    16 & 5 & 277 & 148.21 & 6.71 & 154.92 & 0.41 \\
    17 & 5 & 88 & 11.42 & 3.70 & 15.12 & 0.33 \\
    18 & 6 & 52 & 35.98 & 2.24 & 38.22 & 0.40 \\
    19 & 4 & 49 & 163.19 & 2.67 & 165.86 & 0.26 \\
    20 & 4 & 71 & 8.82 & 3.16 & 11.98 & 0.25 \\
    21 & 6 & 107 & 22.29 & 3.70 & 25.99 & 0.51 \\
    22 & 4 & 62 & 198.65 & 2.42 & 201.07 & 0.18 \\
    23 & 5 & 19 & 115.04 & 1.13 & 116.17 & 0.22 \\
    24 & 4 & 11 & 54.54 & 0.87 & 55.41 & 0.18 \\
    25 & 7 & 33 & 230.96 & 2.05 & 233.01 & 0.24 \\
    26 & 6 & 49 & 265.18 & 1.70 & 266.88 & 0.31 \\
    27 & 11 & 589 & 192.56 & 20.62 & 213.18 & 1.52 \\
    28 & 5 & 206 & 42.73 & 9.34 & 52.07 & 0.38 \\
    29 & 5 & 214 & 88.20 & 6.60 & 94.80 & 0.35 \\
    30 & 6 & 97 & 37.55 & 3.55 & 41.10 & 0.33 \\
    \bottomrule
    \end{tabular}
    }
    \end{table}



\begin{table}[htbp]
    \centering
    \caption{Resultados por PDF: FHM e Hit@k do pipeline textual.}
    \label{tab:resultados-textual-hit-pdf}
    \scriptsize
    \resizebox{\textwidth}{!}{
    \begin{tabular}{r r r r r r r r r}
    \toprule
    \textbf{PDF} & \textbf{Qs} & \textbf{Falhas} & \textbf{FHM} & \textbf{H@1} & \textbf{H@3} & \textbf{H@5} & \textbf{H@10} & \textbf{H@20} \\
    \midrule
    01 & 7 & 0 & 1.7143 & 0.8571 & 0.8571 & 0.8571 & 1.0000 & 1.0000 \\
    02 & 4 & 0 & 1.7500 & 0.5000 & 1.0000 & 1.0000 & 1.0000 & 1.0000 \\
    03 & 7 & 0 & 1.0000 & 0.4286 & 0.4286 & 0.4286 & 0.4286 & 0.4286 \\
    04 & 7 & 0 & 9.6667 & 0.2857 & 0.2857 & 0.5714 & 0.5714 & 0.7143 \\
    05 & 6 & 0 & 15.0000 & 0.0000 & 0.0000 & 0.0000 & 0.1667 & 0.3333 \\
    06 & 7 & 0 & 4.0000 & 0.0000 & 0.4286 & 0.5714 & 0.7143 & 0.7143 \\
    07 & 7 & 0 & 4.4000 & 0.1429 & 0.4286 & 0.4286 & 0.7143 & 0.7143 \\
    08 & 8 & 0 & 2.3750 & 0.7500 & 0.8750 & 0.8750 & 1.0000 & 1.0000 \\
    09 & 6 & 0 & 3.5000 & 0.6667 & 0.8333 & 0.8333 & 0.8333 & 1.0000 \\
    10 & 6 & 0 & 5.3333 & 0.5000 & 0.6667 & 0.6667 & 0.8333 & 1.0000 \\
    11 & 6 & 0 & 2.8333 & 0.3333 & 0.6667 & 0.8333 & 1.0000 & 1.0000 \\
    12 & 5 & 0 & 6.8000 & 0.6000 & 0.8000 & 0.8000 & 0.8000 & 0.8000 \\
    13 & 4 & 0 & 2.0000 & 0.7500 & 0.7500 & 1.0000 & 1.0000 & 1.0000 \\
    14 & 4 & 0 & 1.0000 & 1.0000 & 1.0000 & 1.0000 & 1.0000 & 1.0000 \\
    15 & 5 & 0 & 4.4000 & 0.2000 & 0.6000 & 0.8000 & 0.8000 & 1.0000 \\
    16 & 5 & 0 & 3.2000 & 0.2000 & 0.6000 & 0.8000 & 1.0000 & 1.0000 \\
    17 & 5 & 0 & 9.2000 & 0.4000 & 0.4000 & 0.4000 & 0.4000 & 1.0000 \\
    18 & 6 & 0 & 1.8333 & 0.5000 & 1.0000 & 1.0000 & 1.0000 & 1.0000 \\
    19 & 4 & 0 & 2.2500 & 0.7500 & 0.7500 & 0.7500 & 1.0000 & 1.0000 \\
    20 & 4 & 0 & 2.7500 & 0.7500 & 0.7500 & 0.7500 & 1.0000 & 1.0000 \\
    21 & 6 & 0 & 9.1667 & 0.3333 & 0.3333 & 0.5000 & 0.8333 & 0.8333 \\
    22 & 4 & 0 & 4.3333 & 0.2500 & 0.5000 & 0.5000 & 0.7500 & 0.7500 \\
    23 & 5 & 0 & 5.6000 & 0.4000 & 0.4000 & 0.6000 & 0.8000 & 1.0000 \\
    24 & 4 & 0 & 1.7500 & 0.7500 & 0.7500 & 1.0000 & 1.0000 & 1.0000 \\
    25 & 7 & 0 & 1.0000 & 0.4286 & 0.4286 & 0.4286 & 0.4286 & 0.4286 \\
    26 & 6 & 0 & 19.0000 & 0.1667 & 0.3333 & 0.3333 & 0.5000 & 0.6667 \\
    27 & 11 & 0 & 27.3636 & 0.0909 & 0.0909 & 0.1818 & 0.3636 & 0.4545 \\
    28 & 5 & 0 & 5.7500 & 0.4000 & 0.4000 & 0.4000 & 0.6000 & 0.8000 \\
    29 & 5 & 0 & 3.2000 & 0.4000 & 0.4000 & 1.0000 & 1.0000 & 1.0000 \\
    30 & 6 & 0 & 3.7500 & 0.3333 & 0.5000 & 0.5000 & 0.5000 & 0.6667 \\
    \bottomrule
    \end{tabular}
    }
    \end{table}
    
    
    
    
    \begin{table}[htbp]
    \centering
    \caption{Resultados por PDF: MRR e Evidence Page Recall@k do pipeline textual.}
    \label{tab:resultados-textual-recall-pdf}
    \scriptsize
    \resizebox{\textwidth}{!}{
    \begin{tabular}{r r r r r r r r r}
    \toprule
    \textbf{PDF} & \textbf{Qs} & \textbf{Falhas} & \textbf{MRR} & \textbf{R@1} & \textbf{R@3} & \textbf{R@5} & \textbf{R@10} & \textbf{R@20} \\
    \midrule
    01 & 7 & 0 & 0.8810 & 0.5714 & 0.6429 & 0.7143 & 0.9048 & 0.9524 \\
    02 & 4 & 0 & 0.7083 & 0.2500 & 0.5714 & 0.7321 & 0.7679 & 1.0000 \\
    03 & 7 & 0 & 0.4286 & 0.4286 & 0.4286 & 0.4286 & 0.4286 & 0.4286 \\
    04 & 7 & 0 & 0.3599 & 0.2857 & 0.2857 & 0.3179 & 0.3321 & 0.5393 \\
    05 & 6 & 0 & 0.0380 & 0.0000 & 0.0000 & 0.0000 & 0.1667 & 0.2222 \\
    06 & 7 & 0 & 0.2421 & 0.0000 & 0.1857 & 0.2110 & 0.4330 & 0.5297 \\
    07 & 7 & 0 & 0.3204 & 0.1071 & 0.3214 & 0.3571 & 0.4464 & 0.5000 \\
    08 & 8 & 0 & 0.8042 & 0.5625 & 0.7500 & 0.8125 & 0.8750 & 1.0000 \\
    09 & 6 & 0 & 0.7611 & 0.5833 & 0.6667 & 0.6667 & 0.8333 & 1.0000 \\
    10 & 6 & 0 & 0.5839 & 0.2222 & 0.5000 & 0.5833 & 0.7500 & 0.9167 \\
    11 & 6 & 0 & 0.5655 & 0.2500 & 0.5833 & 0.7500 & 0.9167 & 0.9167 \\
    12 & 5 & 0 & 0.7069 & 0.3667 & 0.5667 & 0.5667 & 0.5667 & 0.6667 \\
    13 & 4 & 0 & 0.8000 & 0.6250 & 0.7500 & 1.0000 & 1.0000 & 1.0000 \\
    14 & 4 & 0 & 1.0000 & 1.0000 & 1.0000 & 1.0000 & 1.0000 & 1.0000 \\
    15 & 5 & 0 & 0.4567 & 0.2000 & 0.5000 & 0.7000 & 0.8000 & 0.9333 \\
    16 & 5 & 0 & 0.4500 & 0.2000 & 0.6000 & 0.7000 & 0.9000 & 0.9500 \\
    17 & 5 & 0 & 0.4418 & 0.3000 & 0.3000 & 0.3000 & 0.3000 & 0.6333 \\
    18 & 6 & 0 & 0.6944 & 0.5000 & 0.7917 & 0.9167 & 0.9583 & 0.9583 \\
    19 & 4 & 0 & 0.7917 & 0.7500 & 0.7500 & 0.7500 & 1.0000 & 1.0000 \\
    20 & 4 & 0 & 0.7812 & 0.5357 & 0.5714 & 0.6071 & 0.8036 & 0.8393 \\
    21 & 6 & 0 & 0.4097 & 0.1667 & 0.3333 & 0.5000 & 0.6944 & 0.6944 \\
    22 & 4 & 0 & 0.3611 & 0.2500 & 0.5000 & 0.5000 & 0.7500 & 0.7500 \\
    23 & 5 & 0 & 0.4782 & 0.4000 & 0.4000 & 0.5000 & 0.7000 & 0.8833 \\
    24 & 4 & 0 & 0.8125 & 0.7500 & 0.7500 & 1.0000 & 1.0000 & 1.0000 \\
    25 & 7 & 0 & 0.4286 & 0.3333 & 0.3810 & 0.3810 & 0.3810 & 0.3810 \\
    26 & 6 & 0 & 0.2632 & 0.0833 & 0.3333 & 0.3333 & 0.5000 & 0.5667 \\
    27 & 11 & 0 & 0.1519 & 0.0909 & 0.0909 & 0.1818 & 0.3182 & 0.3864 \\
    28 & 5 & 0 & 0.4429 & 0.3000 & 0.3000 & 0.3000 & 0.4000 & 0.6000 \\
    29 & 5 & 0 & 0.5300 & 0.3000 & 0.3000 & 0.7667 & 0.9333 & 1.0000 \\
    30 & 6 & 0 & 0.4318 & 0.1389 & 0.3889 & 0.3889 & 0.4444 & 0.5000 \\
    \bottomrule
    \end{tabular}
    }
    \end{table}
    



\subsection{Análise dos casos de baixo desempenho}

Além dos resultados agregados por documento, também foram analisados casos específicos em que o pipeline textual apresentou evidente baixo desempenho. Para essa análise, o ranking foi interpretado no nível dos chunks recuperados: \texttt{hit\_at\_20 = false} indica que nenhuma página de evidência foi encontrada entre as páginas associadas aos 20 primeiros chunks, enquanto \texttt{first\_hit\_rank >= 10} indica que a primeira página correta apareceu apenas a partir do décimo chunk recuperado. Como esses critérios podem se sobrepor, a Tabela~\ref{tab:baixo-desempenho-resumo} apresenta categorias mutuamente exclusivas para resumir os casos observados.

\begin{table}[htbp]
\centering
\caption{Resumo dos casos de baixo desempenho do pipeline textual.}
\label{tab:baixo-desempenho-resumo}
\small
\begin{tabular}{p{9.2cm}r}
\toprule
\textbf{Critério analisado} & \textbf{Quantidade} \\
\midrule
Total de perguntas avaliadas & 172 \\
Sem nenhum acerto entre os 100 chunks recuperados & 20 \\
Primeiro acerto apenas depois do rank 20 & 12 \\
Primeiro acerto entre os ranks 10 e 20 & 17 \\
\bottomrule
\end{tabular}
\end{table}

Além disso, é interessante utilizar o argumento relacionado aos tipos de entidade que acarretaram tais falhas no pipeline de recuperação, visto que é importante objeto de estudo da multimodalidade. Portanto, abaixo seguem os tipos de evidência mais recorrentes em falhas.


\begin{table}[htbp]
\centering
\caption{Falhas do pipeline textual por tipo de evidência.}
\label{tab:falhas-tipo-evidencia}
\small
\begin{tabular}{lrrr}
\toprule
\textbf{Tipo de evidência} & \textbf{Total} & \textbf{Sem H@20} & \textbf{Taxa de falha} \\
\midrule
Figure & 57 & 16 & 28,1\% \\
Chart & 44 & 8 & 18,2\% \\
Table & 34 & 4 & 11,8\% \\
Generalized-text/Layout & 38 & 3 & 7,9\% \\
Pure-text & 40 & 1 & 2,5\% \\
\bottomrule
\end{tabular}
\end{table}

Os resultados indicam que o principal padrão de falha do pipeline textual ocorre em perguntas cuja resposta depende de elementos visuais, como figuras, gráficos, mapas, slides, cores, diagramas ou relações espaciais de layout. Esse comportamento é esperado para uma abordagem baseada apenas em texto extraído, uma vez que informações visuais podem não ser convertidas adequadamente para uma representação textual recuperável.

Observa-se que as maiores taxas de falha ocorreram em evidências do tipo Figure e Chart, com 28,1\% e 18,2\% de falhas, respectivamente. Em contrapartida, é fato de que evidências em texto apresentam taxa de acerto excelente.

A Tabela~\ref{tab:pdfs-problematicos-textual} apresenta os PDFs com maior número de perguntas sem acerto até H@20.

\begin{table}[htbp]
\centering
\caption{PDFs mais problemáticos para o pipeline textual.}
\label{tab:pdfs-problematicos-textual}
\small
\begin{tabular}{r r r r r}
\toprule
\textbf{PDF} & \textbf{Sem H@20} & \textbf{Perguntas} & \textbf{MRR} & \textbf{R@20} \\
\midrule
27 & 6 & 11 & 0.1519 & 0.3864 \\
05 & 4 & 6 & 0.0380 & 0.2222 \\
03 & 4 & 7 & 0.4286 & 0.4286 \\
25 & 4 & 7 & 0.4286 & 0.3810 \\
04 & 2 & 7 & 0.3599 & 0.5393 \\
06 & 2 & 7 & 0.2421 & 0.5297 \\
07 & 2 & 7 & 0.3204 & 0.5000 \\
26 & 2 & 6 & 0.2632 & 0.5667 \\
30 & 2 & 6 & 0.4318 & 0.5000 \\
\bottomrule
\end{tabular}
\end{table}

Analisando as métricas, alguns exemplos ajudam a ilustrar o comportamento observado. No PDF 3, há uma pergunta sobre quanto tempo era gasto com família e amigos em 2010, onde a evidência se encontra em um gráfico localizado na página 14, conforme é observado na~\ref{fig:erros-textual}. Entretanto, as páginas mais bem ranqueadas foram 1, 2, 3, 8 e 11. O motivo da falha é que a informação necessária estava codificada no gráfico demonstrado na imagem, assim dificultando a extração da mesma.

No PDF 6, uma pergunta sobre a comparação entre o índice de ``IPOs'' da Europa e dos Estados Unidos também dependia de um gráfico localizado na página 11, conforme demonstrado pela~\ref{fig:erros-textual}. Nesse caso, as páginas inicialmente recuperadas foram 7, 28, 1, 9 e 30. O pipeline textual recuperou páginas semanticamente próximas, mas não a página visual correta para a consulta.

No PDF 7, uma pergunta sobre o fluxograma ``Analytics Value Chain'' tinha como evidência a página 13. As páginas recuperadas no topo do ranking foram 52, 11, 53, 14 e 15, sugerindo dificuldade do pipeline textual em representar adequadamente a estrutura visual do fluxograma.

Outro exemplo ocorre no PDF 30, em uma pergunta sobre as portas localizadas no lado direito do MacBook Air. A evidência estava associada a uma figura na página 30, mas as páginas mais bem ranqueadas foram 5, 60, 61, 36 e 31.

De modo geral, a análise dos casos de baixo desempenho mostra que o baseline textual falha principalmente quando a pergunta exige interpretação visual, como leitura de gráficos, mapas, cores, diagramas, cartoons, posição espacial ou elementos de layout. Quando a evidência está representada como texto puro, a taxa de falha é significativamente menor. Esse resultado reforça a justificativa para comparar o baseline textual com estratégias visuais e híbridas, como ColQwen e RAG-Anything, que têm acesso mais direto à informação visual presente nos documentos.

\begin{figure}[h!]
    \centering

    \begin{minipage}{0.50\textwidth}
        \centering
        \includegraphics[width=\textwidth]
        {02_Textuais/desenvolvimento/metodologia/imagens/figura_pdf3.png}
        \caption*{(a) Exemplo da página 3}
    \end{minipage}
    \hfill
    \begin{minipage}{0.49\textwidth}
        \centering
        \includegraphics[width=\textwidth]
        {02_Textuais/desenvolvimento/metodologia/imagens/graphpdf6.png}
        \caption*{(b) Exemplo da página 6}
    \end{minipage}

    \caption{Exemplos de páginas associadas a casos de baixo desempenho do pipeline textual.}
    \label{fig:erros-textual}
\end{figure}
